\documentclass[12pt]{article}

% House style preamble
\input{.house-style/preamble.tex}

% Update PDF metadata
\hypersetup{
  pdftitle={What do we mean by language}
}

\title{What do we mean by \mention{language}? \\
  \large A pluralist map for the language sciences}
\author{Brett Reynolds \orcidlink{0000-0003-0073-7195}%
\thanks{Contact: \href{mailto:brett.reynolds@humber.ca}{brett.reynolds@humber.ca}}\\
Humber Polytechnic \& University of Toronto}
\date{\today}

\begin{document}
\maketitle

\begin{abstract}
TODO: Write abstract here.
\end{abstract}

\section{A slippage and a question}

In April 2026, \textcite{casten2026} reported that a small fraction of the human genome carries an outsized share of variation in spoken language ability. Their target measure, the F1 factor, is anchored by sentence repetition and aggregates morphosyntactic, phonological, working-memory, lexical, and motor contributions to verbal performance. In a community-based cohort of 350 children assessed longitudinally from kindergarten through fourth grade, a polygenic score restricted to human-ancestor quickly-evolved regions (HAQERs, less than 0.1\% of the genome) showed an effect of $\beta = 0.18$ on F1, no effect on nonverbal IQ, and replication across more than 100,000 participants in three further cohorts with phenotypically coarser measures. Cross-species analysis of 170 nonprimate mammalian genomes found that HAQER-like sequence similarity correlates with vocal-learning ability across 49 vocal-learning species and 121 non-vocal learners.

Within days of publication, Steven Pinker cited the paper on social media as vindication of \textcite{pinkerbloom1990}, the foundational argument that the human language faculty is a complex adaptation shaped by natural selection for communication. The reading was widely circulated and approvingly relayed. It is also a slippage.

What \citeauthor{casten2026} measured is a developmental verbal-performance phenotype anchored by sentence repetition: a score that loads on attention, working memory, and lexical access alongside grammatical processing. The cross-mammalian regulatory architecture they identify is shared across species that learn vocalizations, including dolphins, parrots, and pinnipeds. Neither finding bears directly on the architecture of a uniquely human grammatical capacity. Pinker--Bloom's faculty -- the modular, recursive, structurally specified system that distinguishes human language from animal communication -- is at a different grain. The shared label \mention{language} obscures the move. A developmental verbal-performance phenotype gets read as evidence about the language faculty because both are referred to as evidence about \mention{language}.

This is not a one-off. The word \mention{language} does work for several recurrent senses across linguistics, cognitive science, anthropology, philosophy, semiotics, law, and education. It picks out an individual mental system in one paragraph and a community-level convention in the next; a developmental phenotype here and an institutionally-stabilized national object there; a set of formal strings in semantic theory and a situated communicative practice in linguistic anthropology. Each of these senses is real, well-studied, and distinct in its projection target. None licenses inferences about the others without further argument. Treating evidence about one as evidence about another is the recurrent pattern that the Casten / Pinker case exemplifies in miniature.

The question of this paper, then: how should we use the word \mention{language} across fields, given that the word does work for senses that don't all license the same inferences? Two answers already on offer fall short. The first is monism: pick one sense and treat the rest as misapplications. \textcite[26]{Chomsky1986} adopts this stance for psychological language and dismisses social senses as \enquote{artificial, somewhat arbitrary, and perhaps not very interesting constructs}. \textcite{santana2016language} shows the dismissal fails: positive arguments for each sense succeed, and the negative arguments for ruling out the others generally fail. The second is unconstrained pluralism: accept the senses as a list and leave it at that. Santana himself rules this out -- the senses interact, and confirmation, idealization, and theory choice all turn on which sense a researcher is committed to.

This paper takes a third route. The senses cluster: they are recurrent, recognisable, and rough enough at the edges that working scientists already treat them as different. Each cluster is real \emph{for the field purposes its mechanisms underwrite} and unreal for purposes those mechanisms cannot project to. Boyd's homeostatic-property-cluster framework~-- categories as profiles of co-occurring properties, stabilized by mechanisms, projectible relative to purposes \citep{boyd1991,boyd1999}~-- supplies the discipline. The contribution is the synthesis: a more granular cluster taxonomy than Santana's three ontologies, plus the diagnostic that lets us say when claims about one cluster smuggle in claims about another. The Casten / Pinker case is a worked instance of exactly that diagnosis. Section~\ref{sec:slippage} works through it in detail.

The map proposed has six clusters: capacity, idiolect, population-level convention, named sociohistorical object, abstract formal system, and communicative practice. The taxonomy is open: new clusters are admissible if they pass the projectibility and homeostasis tests. The umbrella label \mention{language} remains useful as a coordinative term across fields. The clusters under it are what working science actually projects from.

\section{Pluralism, and what HPC adds}

TODO: Engage Santana 2016. Introduce HPC discipline. Set aside Language of thought; note the broader semiotic genus; anticipate operationalization slippage.

\section{Capacity}

TODO: Five sub-axes: faculty / processing-architecture; broad-faculty; developmental-acquisition; comparative-biology / vocal-learning; impairment-phenotype.

\section{Idiolect and \mention{I-language}}

TODO: Three sub-axes: psychological conception (Chomsky); linguistic conception (Devitt); empirical individual variation.

\section{Population-level convention}

TODO: Three sub-axes: coordination convention (Lewis); speech-community variables (Labov, Tagliamonte); stylistic / indexical practice (Eckert).

\section{Named sociohistorical objects}

TODO: Three sub-axes: institutional artefact; identity-marker; descended kind.

\section{Abstract formal systems}

TODO: Four sub-axes: formal-language theory; formal semantics; mathematical-Platonist linguistics; computational / NLP working usage.

\section{Communicative practice}

TODO: Four primary sub-axes plus a cross-cutting modality axis: speech-acts; interactional / situational organization; communicative competence; metapragmatic typification / enregisterment; modality.

\section{Slippage}\label{sec:slippage}

TODO: The methodological payoff. Catalogue slippage patterns. Casten / Pinker as worked example. Operationalization slippages.

\section{The umbrella label}

TODO: \mention{Language simpliciter} as a coordinative label across fields, not a realist kind.

\section{Conclusion}

TODO: Replace the question \enquote{What is language?} with better questions about cluster, projection target, slippage risk.

\section*{Acknowledgements}

% TODO: Replace with the actual models used on this project and their current versions.
% Check model names: frontier models change frequently (e.g., GPT-4 → GPT-4o → o3).
This paper was drafted with the assistance of large language models. All content has been reviewed and revised by the author, who takes full responsibility for the final text.

\printbibliography

\end{document}
